% =====================================================================
% Discussion, Conclusion, back matter, references and appendices,
% revised 2026-10-03. Replaces everything from \section{Discussion} to
% \end{document} in main.tex. New references: refs_methods_additions_20261003.bib
% New figures (R): D1_closure_response, D2_zscores, D3_noise_boxplots
% (scripts/R/make_discussion_figures.R) and M4_parity (moved here).
% =====================================================================
\FloatBarrier
\section{Discussion}
\label{sec:discussion}

\subsection{Principal findings}
This study calibrated a closed-loop lumped-parameter model to the pre-closure catheterization record of one child with a restrictive perimembranous VSD and used it, without refitting, to predict the circulation of the same child after transcatheter closure. The main findings are:
\begin{enumerate}[leftmargin=1.5em,itemsep=1pt]
\item Seven of 14 candidate parameters controlled the nine pre-closure calibration targets. For seven of nine outputs the dominant parameter was the determinant that clinical physiology assigns to that output, and for the other two it agreed in part (Table~\ref{tab:physio_gsa}).
\item Calibration reduced the RMSE of the nine targets from 0.274 to 0.105, and the catheter measurements that were not targets were reproduced within 11\%.
\item The simulated shunt behaved as a restrictive defect. It was left-to-right throughout the cycle, driven by the LV--RV pressure gradient and insensitive to pulmonary resistance.
\item At closure the model reproduced the direction of change of pulmonary flow, $Q_p/Q_s$, PA pressure, mean arterial pressure and RA pressure, but not of systolic and diastolic RFA pressure, and it did not predict the expected unloading of the left heart (Figure~\ref{fig:closure}).
\item The prediction errors fall into three groups: errors within the disagreement between measurement sites and methods (diastolic systemic pressure, $Q_s$), an error carried over from calibration (systolic PA pressure), and errors in chamber volumes that no pre-closure measurement constrains (Figure~\ref{fig:zscores}).
\end{enumerate}

\subsection{Consistency of the sensitivity structure with clinical physiology}
\label{sec:disc_gsa}
Table~\ref{tab:physio_gsa} sets the dominant parameters found by the screen beside the determinant that clinical physiology assigns to each output.

\begin{table}[!htbp]
\centering\small
\setlength{\tabcolsep}{3pt}
\caption{Dominant parameters of the sensitivity screen at the scaled baseline (largest total-order indices, Figure~\ref{fig:gsa}) against the established physiological determinant of each output.}
\label{tab:physio_gsa}
\footnotesize
\begin{tabularx}{\linewidth}{>{\raggedright\arraybackslash}p{2.2cm} >{\raggedright\arraybackslash}X >{\raggedright\arraybackslash}X >{\raggedright\arraybackslash}p{2.2cm}}
\toprule
\textbf{Output} & \textbf{Dominant parameters ($S_T$)} & \textbf{Established determinant} & \textbf{Agreement} \\
\midrule
$Q_p/Q_s$ & $C_d$ 0.51, systemic resistance 0.41, pulmonary resistance 0.01 & restrictive VSD: defect size and LV--RV gradient~\cite{Rao2024,Haas_2024} & consistent \\
Mean arterial pressure & systemic resistance 0.58 & $Q_s \times$ SVR & consistent \\
Diastolic RFA pressure & systemic resistance 0.70 & diastolic runoff through SVR & consistent \\
Systolic RFA pressure & systemic resistance 0.43, $C_{\mathrm{PAR}}$ 0.13, $C_{\mathrm{SAR}}$ 0.12 & mean pressure plus stroke volume over arterial compliance~\cite{FloresGeronimo2021} & partly, compliance weaker than expected \\
Mean and diastolic PA pressure & pulmonary resistance 0.30 and 0.40 & PVR, the steady part of RV afterload~\cite{Kempton2025,Gouwy2025} & consistent \\
Systolic PA pressure & $C_{\mathrm{PAR}}$ 0.39, pulmonary resistance 0.19 & PA compliance, the pulsatile part of RV afterload~\cite{Kempton2025,Gouwy2025} & consistent at the baseline \\
Mean RA pressure & $E_{\mathrm{B,RV}}$ 0.48, $C_{\mathrm{PAR}}$ 0.14, $C_{\mathrm{SVEN}}$ 0.11 & RV filling properties and mean systemic filling pressure~\cite{Persichini2022} & consistent \\
$Q_s$ & systemic resistance 0.33, $C_{\mathrm{PAR}}$ 0.16, $C_{\mathrm{SVEN}}$ 0.14 & venous return: mean systemic filling pressure and resistance to venous return~\cite{Persichini2022} & partly, blood volume not screened \\
LVESV, LVEF, RVESV, RVEF & $E_{\mathrm{A,LV}}$ 0.59 and 0.63, $E_{\mathrm{A,RV}}$ 0.71 and 0.76 & contractility relative to arterial load~\cite{Vriz2021} & consistent \\
\bottomrule
\end{tabularx}
\end{table}

\paragraph{The shunt.} In a restrictive defect the orifice limits flow, and the pressure difference between the ventricles, not the ratio of pulmonary to systemic resistance, sets the shunt~\cite{Rao2024,Haas_2024}. The screen reproduced this hierarchy. $Q_p/Q_s$ depended on the discharge coefficient, which scales the effective orifice area, and on systemic resistance, which sets LV systolic pressure and therefore the driving gradient. Pulmonary resistance had almost no effect ($S_T = 0.01$). The orifice law of Eq.~(\ref{eq:orifice}) is the Bernoulli relation that clinicians apply to the Doppler jet velocity to derive the VSD gradient~\cite{Uppu2022}, so the calibrated $C_d$ has a direct clinical counterpart. In a nonrestrictive defect, pulmonary resistance would be expected to dominate $Q_p/Q_s$, so the retained parameter set is specific to the defect type and should be screened again for each patient.

\paragraph{Systemic and pulmonary pressures.} Mean and diastolic systemic pressure followed systemic resistance, as the steady pressure--flow relation requires. Systolic systemic pressure depended more on resistance than on arterial compliance, whereas pulse-pressure physiology gives compliance a larger role~\cite{FloresGeronimo2021}. The narrow sampling range of $C_{\mathrm{SAR}}$ around a value already set from the measured stroke volume and pulse pressure explains part of this difference. In the pulmonary circuit, the screen separated the steady and pulsatile parts of RV afterload as clinical hemodynamics does. Pulmonary resistance set the mean and diastolic PA pressure, and PA compliance set the systolic value. Kempton et al.\ showed that PVR and PA compliance together describe RV afterload~\cite{Kempton2025}, and Gouwy et al.\ showed in children with pulmonary arterial hypertension that PA compliance carries information independent of PVR~\cite{Gouwy2025}. At the calibrated operating point, pulmonary resistance also dominated systolic PA pressure, because calibration lowered $C_{\mathrm{PAR}}$ and with it the absolute range that was sampled.

\paragraph{Filling pressure and systemic flow.} Mean RA pressure depended on RV passive elastance and, to a lesser degree, on systemic venous compliance, which agrees with the view that RA pressure reflects both RV filling and the mean systemic filling pressure~\cite{Persichini2022}. In Guyton's framework, systemic flow equals venous return, which is set by the mean systemic filling pressure, RA pressure and the resistance to venous return, and the stressed blood volume is the main determinant of the filling pressure~\cite{Persichini2022}. Blood volume was not among the screened inputs, so the screen assigns to arterial resistance part of the control of $Q_s$ that physiology assigns to venous capacitance and blood volume. This is a property of the chosen inputs, not a finding about P1.

\paragraph{Chamber function.} Active elastance controlled end-systolic volume and ejection fraction of both ventricles but had $S_T < 0.10$ for every measured pressure and flow. This is the separation described by ventricular--arterial coupling, in which ejection fraction reflects contractility relative to arterial load, while mean pressures and flows are set mainly by the vascular load and venous return~\cite{Vriz2021,Persichini2022}. The pre-closure catheterization therefore contains no quantity that informs contractility. The chamber parameters were not calibrated, and the post-closure chamber volumes are predictions that no measurement constrains. Lechuga et al.\ found the same split in pulmonary hypertension~\cite{Lechuga2026}, and guidelines for mechanistic cardiovascular models recommend restricting calibration to the parameters the data inform~\cite{Colebank2024Guidelines,Colunga2023}.

\subsection{Consistency of the pre-closure hemodynamics with a restrictive VSD}
The calibrated shunt was left-to-right throughout the cycle because LV pressure exceeded RV pressure at every instant, and 40\% of the shunted volume crossed after aortic valve closure (Section~\ref{sec:results_pre}). The LV therefore had no true isovolumic phases (Figure~\ref{fig:pv}), the expected mechanics of a ventricle with a second, low-pressure outlet. The split of shunt volume between systole and relaxation depends on the ventricular relaxation duration, which was not calibrated. Doppler recording of the VSD jet over the cycle would test it directly.

One pre-closure result does not agree with VSD physiology. A left-to-right VSD loads the left atrium and LV, because the shunted blood returns through the pulmonary veins~\cite{Safa2023,Rao2024}. The calibrated model had a larger RV than LV (RVEDV 43.3 against LVEDV 39.8~mL), and the measured post-closure volumes show the opposite relation (30.5 against 36.1~mL). RV passive elastance reached its lower bound because the fit made the RV more compliant to hold mean RA pressure at 5~mmHg, and a compliant RV fills to a large volume. The RV volume is therefore a by-product of fitting RA pressure without a volume measurement.

\subsection{Consistency of the predicted closure response with clinical physiology}
\label{sec:disc_closure}
Table~\ref{tab:physio_closure} and Figure~\ref{fig:closure} compare the simulated change at closure with the change that the physiology of shunt removal predicts and with the measured change.

\begin{table}[!htbp]
\centering\small
\setlength{\tabcolsep}{3pt}
\caption{Response to closure. Change expected from the physiology of removing a left-to-right shunt, change in the model (calibrated pre-closure to predicted post-closure) and measured change.}
\label{tab:physio_closure}
\resizebox{\linewidth}{!}{%
\begin{tabular}{lllll}
\toprule
\textbf{Quantity} & \textbf{Expected after closure} & \textbf{Model, pre $\rightarrow$ post} & \textbf{Measured, pre $\rightarrow$ post} & \textbf{Model agreement} \\
\midrule
$Q_p$ (L/min) & falls to systemic flow & 3.64 $\rightarrow$ 3.17 & 4.09 $\rightarrow$ -- & consistent \\
$Q_p/Q_s$ (--) & 1.0 & 1.16 $\rightarrow$ 1.00 & 1.19 $\rightarrow$ 0.90 & consistent \\
$Q_s$ (L/min) & unchanged or higher & 3.15 $\rightarrow$ 3.17 & 3.42 (Fick) $\rightarrow$ 2.54 (echo) & consistent, measurement is not \\
Mean PA pressure (mmHg) & lower & 15.8 $\rightarrow$ 14.9 & 15 $\rightarrow$ 13 & consistent \\
Systolic PA pressure (mmHg) & lower & 25.1 $\rightarrow$ 23.9 & 20 $\rightarrow$ 17 & consistent \\
Mean arterial pressure (mmHg) & unchanged or higher & 69.6 $\rightarrow$ 70.2 & 71.3 $\rightarrow$ 75.0 & consistent \\
Systolic RFA pressure (mmHg) & unchanged or higher & 88.1 $\rightarrow$ 90.1 & 100 $\rightarrow$ 89 & consistent, measurement is not \\
Diastolic RFA pressure (mmHg) & unchanged or higher & 52.8 $\rightarrow$ 52.0 & 57 $\rightarrow$ 68 & not consistent \\
Mean RA pressure (mmHg) & unchanged & 5.2 $\rightarrow$ 5.4 & 5 $\rightarrow$ 5 & consistent \\
Mean LA pressure (mmHg) & lower & 7.85 $\rightarrow$ 7.92 & not measured & not consistent \\
LVEDV (mL) & lower (LV unloading)~\cite{Yuan2025} & 39.8 $\rightarrow$ 41.1 & -- $\rightarrow$ 36.1 & not consistent \\
\bottomrule
\end{tabular}}
\end{table}

\begin{figure}[!htbp]
\centering
\includegraphics[width=0.75\linewidth]{D1_closure_response.pdf}
\caption{Change at closure, as a percentage of the pre-closure value, in the model (bars) and in the measurements (diamonds). A missing diamond means that the quantity was not measured in both states.}
\label{fig:closure}
\end{figure}

\paragraph{Where the model agrees.} Closing the orifice removed the shunt, pulmonary flow fell by 13\% to equal systemic flow and $Q_p/Q_s$ became 1.00. Mean, systolic and diastolic PA pressure fell in both the model and the measurement, mean arterial pressure rose in both, and RA pressure stayed at about 5~mmHg. For a shunt of this size ($Q_p/Q_s = 1.19$), physiology predicts small hemodynamic changes at closure~\cite{Rao2024}, and the simulated and measured pressure changes were both small. The model changes were smaller than the measured ones for every pressure except RA pressure (Figure~\ref{fig:closure}), which is expected from a model whose vascular properties do not change at closure.

\paragraph{Where the model does not agree.} Closure removes part of the pulmonary venous return to the left heart, and LV dimensions in children normalize after transcatheter VSD closure, most of the change occurring in the first month~\cite{Yuan2025}. The model predicted no left-heart unloading. Mean LA pressure stayed at 7.9~mmHg and LVEDV rose from 39.8 to 41.1~mL. The cause is visible in the flows. Pulmonary venous return per minute fell by 13\%, but heart rate fell at the same time from 119 to 108~bpm, so the return per beat fell only by 4\% while diastole lengthened. With chamber properties fixed at their pre-closure values, the LV filled to a slightly larger volume. The model also predicted a small fall of diastolic RFA pressure where the measurement rose by 11~mmHg. Arterial resistance and compliance are fixed at closure in the model, so a change in vascular tone after the procedure cannot be reproduced.

\paragraph{Where the measurement does not agree with physiology.} Systemic flow fell from 3.42 to 2.54~L/min and $Q_p/Q_s$ fell below 1.0 after complete closure. Closure does not reduce systemic output, and $Q_p/Q_s < 1$ is not possible without a right-to-left pathway, which the procedure did not create. The pre-closure flows come from the Fick method and the post-closure flows from echocardiography. In pediatric catheterization, Fick output computed with an assumed oxygen consumption has a mean bias of $+27.5\%$ against thermodilution or MRI~\cite{Ralston2023}, and Fick and MRI flows agreed poorly in children after heart transplantation~\cite{Schramm2023}. The change of method between the two states is therefore a sufficient explanation for the measured fall in systemic flow and the rise in derived SVR. The measured fall of systolic RFA pressure (100 to 89~mmHg) is also not expected after closure and coincided with a rise of systolic DAo pressure (91 to 104~mmHg), which points to the measurement site rather than to the circulation.

\subsection{Prediction error relative to clinical measurement variability}
\label{sec:disc_variability}
A prediction can only be as close to the measurement as the measurements are to each other. Figure~\ref{fig:zscores} expresses each error in units of the declared measurement uncertainty. Before closure, every calibration target except systolic PA pressure lay within $\lvert z\rvert \le 1.3$. After closure, most quantities stayed within $\lvert z\rvert < 2$. The exceptions were systolic PA pressure ($z = +8.1$), $Q_s$ ($+5.0$), RVEDV ($+4.2$), mean PA pressure ($+2.9$), diastolic RFA pressure ($-2.4$) and $Q_p/Q_s$ ($+2.2$).

\begin{figure}[!htbp]
\centering
\includegraphics[width=0.75\linewidth]{D2_zscores.pdf}
\caption{Standardized residual $z$ of the calibrated pre-closure model (circles) and of the post-closure prediction (triangles). Dark and light bands, $\lvert z\rvert \le 1$ and $\lvert z\rvert \le 2$.}
\label{fig:zscores}
\end{figure}

The largest errors differ in origin. The diastolic RFA error ($-16$~mmHg) is smaller than the spread of the same diastolic pressure across the three post-closure sites, from 55~mmHg by cuff to 72~mmHg in the DAo. In preterm infants, the limits of agreement between cuff and invasive pressure span about $-26$ to $+14$~mmHg~\cite{Abdelmawla2025}, so an error of this size lies within the variability of the clinical record. The $Q_s$ error ($+25\%$) matches the size of the Fick error in children~\cite{Ralston2023} and the change of method between the two states. In contrast, the systolic PA error is several times larger than the spread of the repeated PA readings (at most 1~mmHg) and was already present before closure, so it is a model error. The RVEDV error is a consequence of the unconstrained chamber parameters (Section~\ref{sec:disc_gsa}). The prediction is therefore as close to the measurements as the clinical record allows for systemic pressure and flow, but not for systolic PA pressure and RV volume.

\subsection{Comparison with predictive modeling studies}
Most patient-specific lumped-parameter studies compare the model with the same data used to calibrate it. Bozkurt et al.\ optimized 23 parameters for three children with dilated cardiomyopathy and matched pressures, volumes and cardiac output within 10\%~\cite{Bozkurt2022}. Owashi et al.\ estimated LV pressure in twelve patients with aortic stenosis with total relative errors of 11.9\%~\cite{Owashi2020}, and Lechuga et al.\ reported $R^2$ of 0.78--0.89 for pulmonary pressure waveforms in pulmonary hypertension~\cite{Lechuga2026}. The pre-closure fit here is of the same order, with an RMSE of relative error of 0.105 and a pressure RMSE of 5.2~mmHg (Figure~\ref{fig:parity}).

Few studies test a forward prediction against measurements taken after the intervention. In transcatheter aortic valve replacement, Khodaei et al.\ and Bahadormanesh et al.\ built separate patient-specific models of the pre- and post-procedural states from the measurements of each state~\cite{Khodaei2021,Bahadormanesh2023}. Ooida et al.\ predicted the effect of mitral edge-to-edge repair from pre-procedural models but compared the post-procedural state only qualitatively~\cite{Ooida2024}. Chen et al.\ predicted hemodynamics after virtual surgery in ten infants with borderline left ventricles and compared the predictions with the surgery chosen clinically, not with post-operative measurements~\cite{Chen2026BLV}. Pant et al.\ modeled one child before and after creation of a Potts shunt~\cite{Pant2022}. Song et al.\ noted that post-operative conditions can differ substantially from the pre-operative state, which limits the accuracy of surgical planning that carries pre-operative boundary conditions forward~\cite{Song2025}. Against this background, the present study gives a quantitative forward test. The seven post-closure catheter pressures, none of which entered calibration, were predicted with an RMSE of 6.9~mmHg against 5.2~mmHg for the calibrated pre-closure pressures (Figure~\ref{fig:parity}). Colebank et al.\ simulated VSD of three sizes without calibration to individual children~\cite{Colebank2026}, and Ferrero et al.\ interpreted catheterization data in intracardiac shunts without predicting a post-intervention state~\cite{Ferrero2026Shunts}.

\begin{figure}[!htbp]
\centering
\includegraphics[width=0.55\linewidth]{M4_parity.pdf}
\caption{Model against measured catheter pressure. Filled circles, pre-closure calibration targets. Open circles, pre-closure pressures that were not targets. Triangles, post-closure prediction. Dashed line, identity.}
\label{fig:parity}
\end{figure}

\subsection{Effect of measurement uncertainty}
Figure~\ref{fig:noise} shows the calibrated parameters and the pre- and post-closure errors across the measurement-noise replicates. \pending{Describe the spread once the 20 replicates are complete: which parameters vary most (panel A), whether the pre-closure errors stay small while the post-closure errors spread (panel B), and whether the measured RVEDV and LVEDV fall inside the range of the predicted volumes.} A volume prediction from this record should therefore be reported with its uncertainty, as recommended for mechanistic cardiovascular models~\cite{Colebank2024Guidelines,Kachabi2025}.

\begin{figure}[!htbp]
\centering
\includegraphics[width=\linewidth]{D3_noise_boxplots.pdf}
\caption{Measurement-noise simulation. (A) Each calibrated parameter divided by its value in the reference calibration (log scale). (B) Model minus measurement, as a percentage of the measurement, for the calibrated pre-closure model (blue) and the post-closure prediction (orange). Boxes show the median and interquartile range, and points show the individual replicates. \pending{n = 20 recalibrated replicates.}}
\label{fig:noise}
\end{figure}

\subsection{Clinical relevance}
The calibrated model gives quantities that the catheterization record does not contain: the shunt flow over the cycle and its split between systole and relaxation, LA pressure, and the pressure--volume loops of both ventricles. Its sensitivity structure matches the physiology that clinicians use to reason about a restrictive VSD, so its parameters have a clinical reading. $C_d$ corresponds to the effective defect size behind the Doppler gradient, systemic resistance to LV afterload, and PA compliance to the pulsatile RV load. From pre-closure data alone, the model predicted post-closure systolic and mean arterial pressure within 7\% and RA pressure within 9\%. It did not predict left-heart unloading or chamber volumes and should not be used for those until a pre-closure volume measurement enters calibration and the chamber parameters can adapt after closure. Clinical decisions on closure rest on defect size, shunt magnitude and symptoms~\cite{Rao2024}. A calibrated model adds a prediction of the post-closure state for the individual child, a step towards cardiovascular digital twins~\cite{Niederer2021,Katsoulakis2024}.

\subsection{Limitations and future work}
\begin{itemize}[leftmargin=1.5em,itemsep=1pt]
\item One patient was studied, so no claim of population-level performance is made.
\item The prediction describes the acute state only. Chamber and vascular parameters are held at their pre-closure values, so the model cannot represent the adaptation of the LV and of vascular tone over the hours and weeks after closure. The study protocol schedules further measurements one day after closure and at the first outpatient visit, and calibrating time-dependent chamber parameters to these data is the next step.
\item No post-closure calibration was performed. Refitting the model to the post-closure record would show which parameters change at closure and would allow what-if simulations of the closed-defect circulation, for example of changes in resistance or compliance.
\item Two of the seven calibrated parameters ended on their lower bounds, and systolic PA pressure remained $5.1\sigma$ above its measurement after calibration. A single pulmonary arterial compartment cannot reproduce a narrow pulse pressure with a normal mean, and a proximal and distal pulmonary arterial compartment would be the structural remedy.
\item No chamber volume was measured before closure, so the chamber parameters were not calibrated and the post-closure volumes are unconstrained predictions.
\item Ventricular activation and relaxation durations come from the adult reference and were not calibrated. The phase split of shunt flow depends on them.
\item Pre- and post-closure systemic flow come from different methods (Fick and echocardiography), and the three post-closure measurement sites disagree on diastolic systemic pressure by 17~mmHg.
\item The sensitivity screen used 128 training simulations stopped at 10 cycles, which supports ranking but not precise indices. It sampled every input independently, without the resistance--compliance coupling used in calibration, and it did not include blood volume.
\end{itemize}

% =====================================================================
\FloatBarrier
\section{Conclusion}
\label{sec:conclusion}
A closed-loop lumped-parameter model of pediatric VSD was scaled by body weight, screened by global sensitivity analysis, calibrated to one child's pre-closure catheterization and used without refitting to predict the same child's post-closure state. Calibration reduced the RMSE of the nine targets from 0.274 to 0.105. After closure, the relative errors were $+1.2\%$ for systolic and $-6.4\%$ for mean arterial pressure, $+3.2\%$ for diastolic PA pressure and $+8.2\%$ for RA pressure. The largest errors were $+40.7\%$ for systolic PA pressure, an error carried over from calibration, and $+42.1\%$ for RVEDV, which no pre-closure measurement constrains. \pending{One sentence on the measurement-noise result.} Adding a pre-closure echocardiographic volume measurement to calibration, and follow-up measurements after closure, are the next steps this work identifies.

% =====================================================================
\section*{Acknowledgments}
\pending{Clinical staff of RSAB Harapan Kita who acquired the data, if they are not authors.}

\section*{Funding}
\pending{Funding source and grant number, or ``This research received no external funding.''}

\section*{Conflict of interest}
\pending{``The authors declare no conflict of interest.'' or the declared conflicts.}

\section*{Author contributions}
\pending{CRediT roles for each author.}

\section*{Data availability}
The de-identified clinical record of P1 and the exported model outputs are available from the corresponding author on reasonable request, subject to the approval of the Institutional Review Board of RSAB Harapan Kita.

\section*{Code availability}
The MATLAB and R code that reproduces every result, table and figure is available at \url{https://github.com/HafizAqilla/unified-vsd}. \texttt{run\_all.m} reproduces the pre-closure pipeline and the post-closure prediction, and \texttt{run\_noise\_simulation.m} reproduces the measurement-noise simulation.

\clearpage
\printbibliography

% =====================================================================
\clearpage
\appendix
\section{Circulation equations}
\label{app:circuit}
Chamber pressures follow Eq.~(\ref{eq:elastance}). Vascular pressures are $P_i = (V_i - V_{0,i})/C_i$ for $i \in \{\mathrm{SAR}, \mathrm{SC}, \mathrm{SVEN}, \mathrm{PAR}, \mathrm{PVEN}\}$, with RA, LA, systemic venous and pulmonary venous pressure bounded below at $-5$~mmHg. Valve flows follow Eq.~(\ref{eq:valve}) and the shunt Eq.~(\ref{eq:orifice}). The systemic capillary resistance is split into two halves around the capillary compliance. The vascular balances are
\begin{align}
\dot{V}_{\mathrm{SAR}} &= Q_{\mathrm{AV}} - Q_{\mathrm{SAR}}, \\
L_{\mathrm{SAR}}\dot{Q}_{\mathrm{SAR}} &= P_{\mathrm{SAR}} - P_{\mathrm{SC}} - \left(R_{\mathrm{SAR}} + \tfrac{1}{2}R_{\mathrm{SC}}\right)Q_{\mathrm{SAR}}, \\
\dot{V}_{\mathrm{SC}} &= Q_{\mathrm{SAR}} - Q_{\mathrm{SC}}, \qquad Q_{\mathrm{SC}} = \frac{P_{\mathrm{SC}} - P_{\mathrm{SVEN}}}{\tfrac{1}{2}R_{\mathrm{SC}}}, \\
\dot{V}_{\mathrm{SVEN}} &= Q_{\mathrm{SC}} - Q_{\mathrm{SVEN}}, \\
L_{\mathrm{SVEN}}\dot{Q}_{\mathrm{SVEN}} &= P_{\mathrm{SVEN}} - P_{\mathrm{RA}} - R_{\mathrm{SVEN}}Q_{\mathrm{SVEN}}, \\
\dot{V}_{\mathrm{PAR}} &= Q_{\mathrm{PV}} - Q_{\mathrm{PAR}}, \\
L_{\mathrm{PAR}}\dot{Q}_{\mathrm{PAR}} &= P_{\mathrm{PAR}} - P_{\mathrm{PC}} - R_{\mathrm{PAR}}Q_{\mathrm{PAR}}, \\
\left(C_{\mathrm{PCOX}} + C_{\mathrm{PCNO}}\right)\dot{P}_{\mathrm{PC}} &= Q_{\mathrm{PAR}} - Q_{\mathrm{PC}}, \qquad Q_{\mathrm{PC}} = \left(P_{\mathrm{PC}} - P_{\mathrm{PVEN}}\right)\left(\frac{1}{R_{\mathrm{PCOX}}} + \frac{1}{R_{\mathrm{PCNO}}}\right), \\
\dot{V}_{\mathrm{PVEN}} &= Q_{\mathrm{PC}} - Q_{\mathrm{PVEN}}, \\
L_{\mathrm{PVEN}}\dot{Q}_{\mathrm{PVEN}} &= P_{\mathrm{PVEN}} - P_{\mathrm{LA}} - R_{\mathrm{PVEN}}Q_{\mathrm{PVEN}}.
\end{align}

\section{Nominal adult reference parameters}
\label{app:nominal_parameters}
Table~\ref{tab:nominal_parameters} gives the nominal parameter set used before allometric scaling. The values were tuned from the adult sets of Valenti~\cite{Valenti2023}, Zhang et al.~\cite{zhang2019multiscale} and Bozkurt~\cite{Bozkurt2019} as described in Section~\ref{sec:model}. Supplementary Section~S6 gives the three source values of every entry.

\begin{table}[!htbp]
\centering\small
\caption{Nominal adult reference values used before allometric scaling.}
\label{tab:nominal_parameters}
\begin{tabularx}{\textwidth}{>{\raggedright\arraybackslash}p{2.7cm} >{\raggedright\arraybackslash}X >{\raggedright\arraybackslash}p{4.0cm} >{\raggedright\arraybackslash}p{2.7cm}}
\toprule
\textbf{Group} & \textbf{Parameters} & \textbf{Nominal values} & \textbf{Unit} \\
\midrule
Heart rate & HR & 75 & bpm \\
LV elastance & $E_{\mathrm{A,LV}}$, $E_{\mathrm{B,LV}}$ & 3.5, 0.08 & mmHg/mL \\
RV elastance & $E_{\mathrm{A,RV}}$, $E_{\mathrm{B,RV}}$ & 0.5, 0.042 & mmHg/mL \\
LA elastance & $E_{\mathrm{A,LA}}$, $E_{\mathrm{B,LA}}$ & 0.35, 0.20 & mmHg/mL \\
RA elastance & $E_{\mathrm{A,RA}}$, $E_{\mathrm{B,RA}}$ & 0.30, 0.10 & mmHg/mL \\
Chamber $V_0$ & $V_{0,\mathrm{RA}}$, $V_{0,\mathrm{RV}}$, $V_{0,\mathrm{LA}}$, $V_{0,\mathrm{LV}}$ & 3.5385, 8.4067, 2.3085, 3.5385 & mL \\
Systemic resistance & $R_{\mathrm{SAR}}$, $R_{\mathrm{SC}}$, $R_{\mathrm{SVEN}}$ & 0.050, 0.80, 0.050 & mmHg\,s/mL \\
Systemic compliance & $C_{\mathrm{SAR}}$, $C_{\mathrm{SC}}$, $C_{\mathrm{SVEN}}$ & 1.33, 1.00, 30.0 & mL/mmHg \\
Systemic inertance & $L_{\mathrm{SAR}}$, $L_{\mathrm{SVEN}}$ & $1.0\times10^{-5}$, $1.0\times10^{-5}$ & mmHg\,s$^2$/mL \\
Systemic vascular $V_0$ & $V_{0,\mathrm{SAR}}$, $V_{0,\mathrm{SC}}$, $V_{0,\mathrm{SVEN}}$ & 90.9, 49.0, 450.0 & mL \\
Pulmonary resistance & $R_{\mathrm{PAR}}$, $R_{\mathrm{PCOX}}$, $R_{\mathrm{PCNO}}$, $R_{\mathrm{PVEN}}$ & 0.010, 0.0630, 1.2640, 0.020 & mmHg\,s/mL \\
Pulmonary compliance & $C_{\mathrm{PAR}}$, $C_{\mathrm{PCOX}}$, $C_{\mathrm{PCNO}}$, $C_{\mathrm{PVEN}}$ & 5.000, 0.1983, 0.0017, 13.750 & mL/mmHg \\
Pulmonary inertance & $L_{\mathrm{PAR}}$, $L_{\mathrm{PVEN}}$ & $1.0\times10^{-5}$, $1.0\times10^{-5}$ & mmHg\,s$^2$/mL \\
Pulmonary vascular $V_0$ & $V_{0,\mathrm{PAR}}$, $V_{0,\mathrm{PCOX}}$, $V_{0,\mathrm{PCNO}}$, $V_{0,\mathrm{PVEN}}$ & 196.0, 0, 0, 0 & mL \\
Valve resistance & $R_{\mathrm{open}}$, $R_{\mathrm{closed}}$ & $4.0\times10^{-3}$, $9.4168\times10^4$ & mmHg\,s/mL \\
Smoothing widths & $\varepsilon_{\mathrm{v}}$, $\varepsilon_{\mathrm{s}}$ & 0.5, 0.1 & mmHg \\
Shunt & $C_d$ (adult reference), $\rho$ & 0.7, 1060 & --, kg/m$^3$ \\
Ventricular timing & $T_{vc,\mathrm{LV}}$, $T_{vr,\mathrm{LV}}$, $T_{vc,\mathrm{RV}}$, $T_{vr,\mathrm{RV}}$ & 0.265, 0.40, 0.30, 0.40 & fraction of $T_{\mathrm{HB}}$ \\
LA timing & $t_{ac,\mathrm{LA}}$, $T_{ac,\mathrm{LA}}$, $T_{ar,\mathrm{LA}}$ & 0.75, 0.10, 0.80 & fraction of $T_{\mathrm{HB}}$ \\
RA timing & $t_{ac,\mathrm{RA}}$, $T_{ac,\mathrm{RA}}$, $T_{ar,\mathrm{RA}}$ & 0.80, 0.10, 0.70 & fraction of $T_{\mathrm{HB}}$ \\
\bottomrule
\end{tabularx}
\end{table}

\end{document}
